\section {The five card trick (Combinatorics)}

A magician and an assistant perform the following trick. Someone from the audience picks any $5$ cards from a standard deck of $52$ cards
and hands them to the assistant. The assistant looks at them, hides one of them, and puts the other $4$ face up on the table in a row, in an order of his choice.
The magician, who did not see anything until now, looks at the $4$ cards and names the hidden card.
There is no other communication between them. How is the trick done?

\subsection{Solution}

There are $5$ cards and only $4$ suits, so at least two of the cards have the same suit.
One of these two is hidden and the other one is put first, which tells the magician the suit of the hidden card.

Now the value. Arrange the $13$ values in a circle: $A, 2, 3, \ldots, 10, J, Q, K$ and then $A$ again.
Going clockwise from one of the two cards to the other takes some number of steps, and going the other way takes $13$ minus that number.
One of the two ways is at most $6$ steps. The assistant hides the card that is reached at the end of the short way and puts the other card first.
So the magician knows that the hidden card is $1$ to $6$ steps clockwise from the first card.
For example, if the two cards are the $Q$ and the $3$ of hearts, the assistant hides the $3$, which is $4$ steps after the $Q$ ($K, A, 2, 3$).

The remaining $3$ cards encode the number of steps. Agree on an order of the whole deck, for example by value and then by suit,
so that among the $3$ cards there is a low, a middle and a high card. They can be put down in $3! = 6$ orders, one for each number of steps:
\begin{center}
\begin{tabular}{c c}
low, middle, high & 1 \\
low, high, middle & 2 \\
middle, low, high & 3 \\
middle, high, low & 4 \\
high, low, middle & 5 \\
high, middle, low & 6 \\
\end{tabular}
\end{center}
The magician reads the suit and the starting value from the first card, the number of steps from the order of the other three, and names the hidden card.
This trick is due to the mathematician William Fitch Cheney.
